\documentclass[10pt]{amsart}
\usepackage[a4paper,margin=23mm]{geometry}
\usepackage{amsmath,amssymb,amsthm,mathtools,hyperref,graphicx,xspace,letltxmacro,xpatch}
\usepackage[english]{babel}
\newtheorem{theo}{Theorem}[section]
\newtheorem{lemm}[theo]{Lemma}\newtheorem{prop}[theo]{Proposition}\newtheorem{coro}[theo]{Corollary}\newtheorem{conj}[theo]{Conjecture}
\theoremstyle{definition}\newtheorem{defi}[theo]{Definition}\newtheorem{rema}[theo]{Remark}\newtheorem{exam}[theo]{Example}\newtheorem{exem}[theo]{Example}\newtheorem{nota}[theo]{Notation}
\newcommand{\ie}{i.e.}\newcommand{\eg}{e.g.}\newcommand{\resp}{resp.}\newcommand{\skpt}{\par\smallskip\noindent}
\emergencystretch=3em\pagestyle{plain}
\usepackage[all]{xy}
\usepackage{mathtools}
\usepackage{xpatch}
%\DeclareMathOperator*{\colim}{colim}





\newcommand{\noopsort}[1]{}

\newcommand{\cat}[1]{{\mathbf{#1}}}
\renewcommand{\rm}[1]{{\mathrm{#1}}}

\DeclareMathOperator{\opp}{op}
\DeclareMathOperator{\oDD}{odd}

\DeclareMathOperator{\Ass}{Ass}
\DeclareMathOperator{\holim}{holim}
\DeclareMathOperator{\Ext}{Ext}
\DeclareMathOperator{\kerr}{ker}

\DeclareMathOperator{\Hom}{Hom}
\DeclareMathOperator{\mapp}{map}
\DeclareMathOperator{\modd}{mod}
\DeclareMathOperator{\HH}{HH}
\DeclareMathOperator{\THH}{THH}
\DeclareMathOperator{\BG}{BG}
\DeclareMathOperator{\BU}{BU}
\DeclareMathOperator{\SUU}{SU}


\DeclareMathOperator{\ppt}{pt}


\DeclareMathOperator{\Alg}{Alg}
\DeclareMathOperator{\aalg}{alg}

\DeclareMathOperator{\id}{id}
\DeclareMathOperator{\add}{ad}
\DeclareMathOperator{\Der}{Der}

\DeclareMathOperator{\AQ}{AQ}
\DeclareMathOperator{\sym}{Sym}
\DeclareMathOperator{\Gerr}{Ger}

\DeclareMathOperator{\rlie}{rLie}
\DeclareMathOperator{\rTor}{Tor}
\DeclareMathOperator{\rBar}{Bar}
\DeclareMathOperator{\rSq}{Sq}

\newcommand{\alg}{\mathbf{Alg}}
\newcommand{\AssAlg}{\mathbf{AssAlg}}
\newcommand{\grVect}{\mathbf{grVect}}
\newcommand{\bMod}{\mathbf{Mod}}

\newcommand{\bSet}{\mathbf{Set}}
\newcommand{\bAb}{\mathbf{Ab}}
\newcommand{\rAb}{\mathrm{Ab}}

\newcommand{\bfC}{\mathbf{C}}
\newcommand{\bfD}{\mathbf{D}}

\renewcommand{\L}{\mathrm{L}}
\newcommand{\W}{\mathrm{W}}
\newcommand{\ZZ}{\mathbb{Z}}
\newcommand{\QQ}{\mathbb{Q}}
\newcommand{\RR}{\mathbb{R}}
\newcommand{\FF}{\mathbb{F}}
\newcommand{\clG}{\mathcal{G}}
\newcommand{\clD}{\mathcal{D}}
\newcommand{\clC}{\mathcal{C}}

\newcommand{\fkg}{\mathfrak{g}}
\newcommand{\bbL}{\mathbb{L}}
\newcommand{\bbN}{\mathbb{N}}
\newcommand{\bbQ}{\mathbb{Q}}
\newcommand{\bbRP}{\mathbb{RP}}

\renewcommand*{\to}{\mathchoice{\longrightarrow}{\rightarrow}{\rightarrow}{\rightarrow}}

\let\oldmapsto\mapsto
\renewcommand*{\mapsto}{\mathchoice{\longmapsto}{\oldmapsto}{\oldmapsto}{\oldmapsto}}

\newcommand*{\RightArrow}{\mathchoice{\Longrightarrow}{\Rightarrow}{\Rightarrow}{\Rightarrow}}

\SelectTips{cm}{12}


\newcommand*{\mk}{\mkern -1mu}
\newcommand*{\Mk}{\mkern -2mu}
\newcommand*{\mK}{\mkern 1mu}
\newcommand*{\MK}{\mkern 2mu}

%\hypersetup{urlcolor=purple, linkcolor=blue, citecolor=red}

\newcommand*{\relabel}{\renewcommand{\labelenumi}{(\theenumi)}}
\newcommand*{\romanenumi}{\renewcommand*{\theenumi}{\roman{enumi}}\relabel}
\newcommand*{\Romanenumi}{\renewcommand*{\theenumi}{\Roman{enumi}}\relabel}
\newcommand*{\alphenumi}{\renewcommand*{\theenumi}{\alph{enumi}}\relabel}
\newcommand*{\Alphenumi}{\renewcommand*{\theenumi}{\Alph{enumi}}\relabel}
\let\oldtilde\tilde
\renewcommand*{\tilde}[1]{\mathchoice{\widetilde{#1}}{\widetilde{#1}}{\oldtilde{#1}}{\oldtilde{#1}}}
%%%%%%%%%%%%%%%%%%%%%%%%%%%%%%%%%%%%%%%%%%%%%%%%%%%%%%

\begin{document}
\title{Even associative Quillen cohomology for exterior-algebra diagrams with injective odd coefficient diagrams}
\author{Geoffroy Horel}
\maketitle
\subsection*{Conventions}

We denote by $\grVect_k$ the category of graded vector spaces over a field $k$. This category has a symmetric monoidal structure whose symmetry isomorphism involves the usual sign. Our graded vector spaces are cohomologically graded. We write $V\mapsto sV$ for the shift in this category given by $(sV)^i=V^{i-1}$.

Given a graded vector space $V$, we denote by $\sym(V)$ the symmetric algebra on $V$, i.e. the free graded commutative algebra on $V$. If $V$ is concentrated in odd degrees, we write $\Lambda(V)$ for the exterior algebra on $V$ which is, by definition, the free \emph{strictly} commutative algebra on $V$ (strictly commutative means that elements of odd degrees square to zero). Of course, by the sign rules, there is an isomorphism $\Lambda(V)\cong \sym(V)$ if the characteristic of the field is different from $2$.


\setcounter{section}{1}
\section{Quillen cohomology}

\subsection{Non-abelian derived functor}

We follow the treatment of~\cite{franklandbehavior}. Let $\bfC$ be a complete and cocomplete category with a set of projective compact generators. Following Frankland, such a category shall be called quasi-algebraic. Thanks to a theorem of Quillen (see~\cite[Theorem~2.4]{quillenhomotopical}), if $\bfC$ is a quasi-algebraic category, its category of simplicial objects, denoted $s\bfC$ has a model structure transferred along the right adjoint functor
\begin{align*}
s\bfC&\to s\bSet^{\clG}\\
X_\bullet&\mapsto \{\Hom_{\bfC}(G,X_\bullet)\}_{G\,\in\,\clG}
\end{align*}
where $\clG$ is a set of compact projective generators of $\bfC$.

Given a left adjoint functor between two quasi-algebraic categories, we define its left derived functor
\[
\L F: s\bfC\to s\bfD
\]
to be the left derived functor in the sense of model categories of the functor $F\!\!:\! s\bfC\to s\bfD$ given by degreewise application of the functor $F$.

\begin{rema}\label{remark : non-abelian derived}
There is a conceptual description of the $\infty$-category underlying the model category $s\bfC$ as the \emph{non-abelian derived category of} $\bfC$. Explicitly, this is completion of the category $\bfC^{cp}$ of compact projective objects of $\bfC$ under $\infty$-sifted colimits. Using this terminology, the functor $\L F$ is the unique extension of $F_{|\bfC^{cp}}$ into a functor that preserves sifted colimits. This is originally due to Lurie (see~\cite[Section~5.5.8]{luriehigher} or~\cite[Section~5.1.1]{kestupurity} for a concise version of Lurie's work).
\end{rema}


\subsection{Quillen cohomology}

Let $\bfC$ be a quasi-algebraic category. For $c$ an object of $\bfC$, we may consider the category $\bAb(\bfC/c)$ of abelian group objects over $c$. By the adjoint functor theorem, there is an abelianization functor
\[
\rAb: \bfC/c\to\bAb(\bfC/c)
\]
which is left adjoint to the forgetful functor. We define the cotangent complex of $c$ to be the left derived functor of the abelianization functor. We denote this object by $L_c^{\bfC}$ or simply $L_c$ if there is no ambiguity. It follows from~\cite[Propositions~3.33 and 3.34]{franklandbehavior} that the categories $\bfC/c$ and $\rAb(\bfC/c)$ are quasi-algebraic and hence that the relevant model structures exist.

Given an object $m\in\bAb(\bfC/c)$, we define the Quillen cohomology of $c$ with coefficients in $m$ as~:
\[
\AQ^s_{\bfC}(c,m):=\pi^s\Hom_{\bAb(\bfC/c)}(L_c,m).
\]
(Observe that $\Hom_{\bAb(\bfC/c)}(L_c,m)$ is a cosimplicial abelian group, and we denote by $\pi^s$ the $s^{\rm{th}}$ cohomology group of its associated cochain complex.)

\begin{rema}
In all the cases of interest to us, the category $\bfC$ will come with a forgetful functor to the category of graded vector spaces over $k$. Moreover, it will be the case that the shift functor on the category of graded vector space will pass to the category of abelian group objects in $\bfC/c$. In this case, the André--Quillen cohomology groups are bigraded
\[
\AQ^{s,t}_{\bfC}(c,m)=\AQ^s_{\bfC}(c,s^tm).
\]
\end{rema}


\setcounter{section}{3}
\section{Computation of the obstruction group}

\begin{defi}
We call a bigraded abelian group even (resp. odd) if it is concentrated in bidegrees $(s,t)$ such that $s+t$ is even (resp. odd).
\end{defi}


\begin{lemm}\label{lemm : Hochschild}
Let $A=\Lambda(V)$ be the exterior algebra on a graded vector space ~$V$ concentrated in odd degrees and of finite total dimension. Let $M$ be an $A$-module concentrated in odd degrees and finite-dimensional in each degree. Since $A$ is commutative, we may view $M$ as an $A$-bimodule. Then
\[
\AQ_{\Ass\alg}^{*,*}(A,M)
\]
is even.
\end{lemm}

\begin{proof}
There is an isomorphism
\begin{align*}
\AQ^{s,*}(A,M)&\cong \Der(A,M)\qquad\,\text{if }\;s=0\\
&\cong \HH^{s+1}(A,M)\quad\text{if }\;s>0
\end{align*}
Since $\Der(A,M)\subset \Hom_k(A,M)$ is concentrated in even degrees, it suffices to prove that $\HH^{*,*}(A,M)$ is odd.

There exists a unique map of commutative algebras
\[
A\to A\otimes A
\]
sending $v\in V$ to $v\otimes 1-1\otimes v$. Indeed, if $k$ is of characteristic different from $2$, then $\Lambda(V)\cong\sym(V)$ is free and if $k$ is of characteristic $2$, then we do indeed have
\[
(v\otimes 1-1\otimes v)^2=0
\]
By~\cite[Theorem~X.6.1]{cartaneilenberg}, we have an isomorphism
\[
\HH^*(A,M)\cong\Ext^*_A(k,\tilde{M})
\]
where $\tilde{M}$ is $M$ equipped with the $A$-module structure induced by restriction of scalars along the map $A\to A\otimes A$ described above. By definition of the $A$-bimodule structure on $M$, the $A$-module $\tilde{M}$ is then simply the $k$-vector space $M$ with the trivial $A$-module structure. So, using the finite dimensionality of $M$, we have
\[
\Ext_A^*(k,\tilde{M})\cong \Ext^*_A(k,k)\otimes_kM
\]
Since $M$ is odd, it suffices to prove that $\Ext^*_A(k,k)$ is even. If $V$ is one-dimensional so that $A=k[x]/x^2$. Then a free resolution of $k$ as an $A$-module is given by
\[
k\leftarrow{A}\xleftarrow{\times x} s^{|x|}A \xleftarrow{\times x}s^{2|x|}A\xleftarrow{\times x}\ldots
\]
It follows that the bigraded abelian group $\Ext^{*,*}_A(k,k)$ is even. For a general $V$, $\Ext_A(k,k)$ is a tensor product of finitely many bigraded abelian groups of this form, therefore the answer is still even.
\end{proof}
\setcounter{section}{4}
\section{Diagrams of \texorpdfstring{$E_2$}{E\_2}-algebras}

\subsection{Main theorem}

Let $\bfC$ be a quasi-algebraic category and let $I$ be a small category. In this situation $\bfC^I$ is also a quasi-algebraic category. Given $c: I\to\bfC$ and $m\in\bAb(\bfC^I/c)$, we denote by $\AQ_{\bfC,I}(c,m)$ the corresponding Quillen cohomology.

We start with the following proposition for which we could not find a reference.

\begin{prop}
Let $I$ be a small category, let $A$ be an associative algebra in $\grVect^I$ and let $M$ be an $A$-$A$-bimodule in $\grVect^I$. Then
\[
\AQ^*_{\Ass,I}(A,M)\cong \Ext^*_{A\otimes A^{\opp}}(\Omega_A,M)
\]
where $\Omega_A$ is the bimodule of associative differentials defined by the following exact sequence
\[
0\to \Omega_A\to A\otimes A\xrightarrow{m}A
\]
where $m$ denotes the multiplication map.
\end{prop}

\begin{proof}
This proposition is very classical if $I=[0]$. The object $\Omega_A$ is the result of applying the abelianization functor to $\id_A: A\to A$ viewed as an object of $\AssAlg^I/A$. By definition of Quillen cohomology, this proposition boils down to proving an equivalence
\[
\bbL\rAb(\id_A)\xrightarrow{\simeq} \Omega_A=\rAb(\id_A)
\]
where $\rAb: s\AssAlg/A\to s\bMod_{A\otimes A^{\opp}}$ is the abelianization functor. Moreover, it is enough to prove this equivalence pointwise.

For any $i$ in $I$ and any category $\bfC$, let us denote by $ev_i$ the evaluation functor $\bfC^I\to \bfC$. According to~\cite[Theorem~4.7]{franklandbehavior}, there is a commutative diagram of Quillen left adjoints (the theorem applies since $ev_i$ preserves projective objects):
\[
\xymatrix{ s\left(\AssAlg^I/A\right)\ar[d]_{ev_i}\ar[r]^{\rAb}&s\bMod_{A\otimes A^{\opp}}\ar[d]^{ev_i}\\
s(\AssAlg/A(i))\ar[r]_{\rAb}&s\bMod_{A(i)\otimes A(i)^{\opp}} }
\]
Applying this to $\id_A$ and deriving the functors, we obtain a weak equivalence
\[
(\bbL\Omega_A)(i)\simeq \bbL\Omega_{A(i)}
\]
Now, using the case $I=[0]$ of the proposition, we find
\[
(\bbL\Omega_A)(i)\simeq \Omega_{A(i)}=\Omega_A(i)
\]
as desired.
\end{proof}
\begin{prop}
Let $V: I\to\grVect$ be a diagram taking values in finite dimensional graded vector spaces concentrated in odd degrees. Let $M: I\to \grVect$ be a $\Lambda(V)$-module satisfying pointwise the conditions of Lemma~\ref{lemm : Hochschild}. We moreover assume that $M$ viewed as a diagram of $k$-vector spaces is injective. Then $\AQ_{\Ass,I}(\Lambda(V),M)$ is even.
\end{prop}
\begin{proof}
In order to prove this proposition, we shall first construct a spectral sequence computing the groups $\AQ_{\Ass,I}(\Lambda(V),M)$.



For $A$ an associative algebra in $\grVect^I$ and $M$ an $A$-bimodule, using the previous proposition, we have an isomorphism
\[
\AQ^*_{\Ass,I}(A,M)=\pi_{-*}\RR\Hom_{A\otimes A^{\opp}}(\Omega_A,M).
\]

In order to compute this derived Hom, we can use first the bar resolution in the category of bimodules and obtain
\[
\RR\Hom_{A\otimes A^{\opp}}(\Omega_A,M)=\holim_{[n]\,\in\,\Delta}\RR\Hom_{\grVect^I}\left((A\otimes A^{\opp})^{\otimes n}\otimes\Omega_A,M\right)
\]



Now, in general, if $\bfC$ is an $\infty$-category and $I$ a small category, we have
\[
\mapp_{\bfC^I}(F,G)\simeq\holim_{\Delta}\left([s]\mapsto\prod_{i_0\to i_1\to\ldots\to i_s} \mapp(F(i_0),G(i_s))\right).
\]
So we can write the derived Hom as the limit of the following double cosimplicial diagram
\[
([s],[n])\mapsto \left(\prod_{i_0\to i_1\to\ldots\to i_s}\RR\Hom_{\grVect}\bigl((A\otimes A^{\opp})^{\otimes n}\otimes\Omega_A(i_0),M(i_s)\bigr)\right)
\]


Taking the limit in the $[n]$ direction, we get
\[
\RR\Hom(\Omega_A^1,M)=\holim_{[s]\,\in\,\Delta}\left(\prod_{i_0\to i_1\to\ldots\to i_s}\RR\Hom_{A(i_0)\otimes A(i_0)^{\opp}}\left(\Omega_{A(i_0)},M(i_s)\right)\right)
\]
As for any cosimplicial object, we obtain an associated Bousfield--Kan spectral sequence whose $E_1$ page is given by
\[
E_1^{s,t}=\prod_{i_0\to i_1\to\ldots\to i_s}\AQ^{t}_{\Ass}(A(i_0),M(i_s)).
\]
This spectral sequence converges as it can also be viewed as the spectral sequence of a third quadrant double complex (coming from the above bicosimplicial abelian group after taking the associated complex in both cosimplicial direction). From the comparison between Quillen cohomology and Hochschild cohomology, we see that the row $t=0$ is given by the cosimplicial abelian group
\[
[s]\mapsto \prod_{i_0\to i_1\to\ldots\to i_s}\Der(A(i_0),M(i_s)),
\]
while for $t>0$, we get
\[
[s]\mapsto \prod_{i_0\to i_1\to\ldots\to i_s}\HH^{t+1}(A(i_0),M(i_s)).
\]

So far we did not use anything about the specific situation and the above discussion would apply to any pair $(A,M)$. Using the computation of Lemma~\ref{lemm : Hochschild}, we see that the row $t>0$ of the $E_1$-page is of the form
\[
[s]\mapsto \prod_{i_0\to i_1\to\ldots\to i_s}\Hom_k(F(i_0),M(i_s)),
\]
where $F: I\to\grVect_k$ is the degreewise dual of the functor
\[
i\mapsto \Ext_{A(i)}(k,k).
\]
Likewise the $0^{\rm{th}}$ row is simply given by the cosimplicial abelian group
\[
[s]\mapsto \prod_{i_0\to i_1\to\ldots\to i_s}\Hom_k(V(i_0),M(i_s)).
\]
(indeed there is an isomorphism $\Der(\Lambda(V),M)\cong\Hom_k(V,M)$). From this observation, using injectivity of $M$, we deduce that the $E_2$-page of the spectral sequence is concentrated on the column $s=0$ and
\[
E_2^{0,-t}=\ker\left(d_1: \prod_{i\in I}\AQ^{t}(A(i),M(i))\to \prod_{f: i\to j}\AQ^{t}(A(i),M(j))\right).
\]
In particular, we see that
\[
\AQ^{*,*}_I(A,M)\subset\prod_{i\in I}\AQ^{*,*}(A(i),M(i))
\]
and is therefore even by Lemma~\ref{lemm : Hochschild}.
\end{proof}
\begin{thebibliography}{ČS24}
\bibitem[Fra15]{franklandbehavior} Martin Frankland, Behavior of Quillen (co)homology with respect to adjunctions, Homology Homotopy Appl. 17 (2015), no. 1, 67–109.
\bibitem[Qui67]{quillenhomotopical} Daniel G. Quillen, Homotopical algebra, Lecture Notes in Mathematics, vol. 43, Springer, 1967.
\bibitem[Lur09]{luriehigher} Jacob Lurie, Higher topos theory, Annals of Mathematics Studies, vol. 170, Princeton University Press, 2009.
\bibitem[ČS24]{kestupurity} Kęstutis Česnavičius and Peter Scholze, Purity for flat cohomology, Ann. Math. (2) 199 (2024), no. 1, 51–180.
\bibitem[CE56]{cartaneilenberg} Henri Cartan and Samuel Eilenberg, Homological algebra, Princeton Mathematical Series, vol. 19, Princeton University Press, 1956.
\end{thebibliography}
\end{document}
